\documentclass[11pt,a4paper]{article}
\usepackage[utf8]{inputenc}\usepackage[T1]{fontenc}\usepackage[italian,english]{babel}
\usepackage{amsmath,amssymb,amsthm}\usepackage[margin=2.2cm]{geometry}\usepackage{longtable,booktabs,array,calc}
\usepackage{listings,xcolor}\usepackage{hyperref}\usepackage{url}
\providecommand{\tightlist}{\setlength{\itemsep}{0pt}\setlength{\parskip}{0pt}}
\lstset{basicstyle=\ttfamily\scriptsize,breaklines=true,frame=single,captionpos=b,columns=fullflexible,keepspaces=true,inputencoding=utf8,extendedchars=true}
\title{Proof Pursuit --- Problem 4, part 6\\ \large prove, decisioni degli agenti, codice verificato, letteratura}
\author{Team Proof Pursuit (Researcher: T. Tumini; Referee: gabundos; Creative: M. Renzi; gio3r3gio)}
\date{\today}
\begin{document}\maketitle\tableofcontents
\section*{Come leggere questo rapporto}
Per ogni cella: la richiesta ufficiale, lo stato dichiarato dal team (mai ``risolto'' senza revisione umana), la consegna
con la prova, poi la traccia delle decisioni degli agenti (Researcher: famiglia e motivo; Referee: verdetto ed errore
fatale; Creative: quando e perché è intervenuto) e il codice su cui il tentativo poggia con l'esito della riesecuzione
fidata. DIMOSTRATO, CITATO, VERIFICATO SU INTERVALLO FINITO e CONGETTURA restano distinti. L'architettura del sistema
è descritta nel README del repository.
\section{Problem 4 — Disjoint congruence classes and large gcd}
\subsection{Cella 6 (13 punti) — stato: parziale}
\subparagraph{Parte 6 (C6) --- Beyond the
boundary}\label{parte-6-c6-beyond-the-boundary}

\textbf{Punteggio:} 13 points · \textbf{Valutazione:} Judged ·
\textbf{Open question}

Three directions beyond the certified range; any one of them counts. Any
of the following.

\begin{enumerate}
\def\labelenumi{(\alph{enumi})}
\item
  Decide a size \(k \ge 25\).
\item
  The asymptotic form. It is known that a pairwise disjoint family of
  size \(k\) always has a pair with \[
  \gcd(m_i, m_j) \;\ge\; k \cdot \exp\!\left( -(2 + o(1)) \frac{\log k}{\log\log k} \right),
  \] which is \(k^{1 - o(1)}\) but not linear in \(k\). Prove the
  statement in full, or prove the weaker bound \(\gcd(m_i, m_j) \ge ck\)
  for some absolute constant \(c > 0\), or improve the exponential
  factor above.
\item
  The group form. Let \(G\) be a group, let \(G_1, \ldots, G_k\) be
  subgroups of finite index \(n_i = [G : G_i]\), and let
  \(x_1 G_1, \ldots, x_k G_k\) be pairwise disjoint cosets. Is there a
  pair \(i < j\) with \(\gcd(n_i, n_j) \ge k\)? This is known for
  \(k \le 5\) and open for every \(k \ge 6\); settling \(k = 6\) counts
  as progress.
\end{enumerate}

\emph{(Fine del problema 4. Punteggi: 1+2+3+5+8+13 = 32.)}

\textbf{Enunciato protetto dato agli agenti (state.json).} \texttt{OPEN. Any of: (a) decide a size k \textgreater{}= 25; (b) prove gcd \textgreater{}= ck for an absolute c \textgreater{} 0 or improve the known factor k exp(-(2+o(1)) log k / log log k); (c) the group form (disjoint cosets of finite-index subgroups), settle k = 6.}

\textbf{Evidenza registrata in STATUS.md.} bozza parziale in inglese in submission/ (parziali.py)

\subsubsection{Consegna (prova)}
\paragraph{Problema 4 --- Parte 6 --- bozza di consegna
PARZIALE}\label{problema-4-parte-6-bozza-di-consegna-parziale}

\textbf{Stato dichiarato: PARZIALE.} Nessuna soluzione completa: qui
sotto ciò che è stabilito, la formalizzazione, la posizione rispetto
alla letteratura e ciò che resta aperto. Nulla è dichiarato dimostrato
oltre quanto scritto.

\subparagraph{1. Risultato e ambito}\label{risultato-e-ambito}

\textbf{Richiesta ufficiale.} \#\# Parte 6 (C6) --- Beyond the boundary

\textbf{Punteggio:} 13 points · \textbf{Valutazione:} Judged ·
\textbf{Open question}

Three directions beyond the certified range; any one of them counts. Any
of the following.

\begin{enumerate}
\def\labelenumi{(\alph{enumi})}
\item
  Decide a size \(k \ge 25\).
\item
  The asymptotic form. It is known that a pairwise disjoint family of
  size \(k\) always has a pair with \[
  \gcd(m_i, m_j) \;\ge\; k \cdot \exp\!\left( -(2 + o(1)) \frac{\log k}{\log\log k} \right),
  \] which is \(k^{1 - o(1)}\) but not linear in \(k\). Prove the
  statement in full, or prove the weaker bound \(\gcd(m_i, m_j) \ge ck\)
  for some absolute constant \(c > 0\), or improve the exponential
  factor above.
\item
  The group form. Let \(G\) be a group, let \(G_1, \ldots, G_k\) be
  subgroups of finite index \(n_i = [G : G_i]\), and let
  \(x_1 G_1, \ldots, x_k G_k\) be pairwise disjoint cosets. Is there a
  pair \(i < j\) with \(\gcd(n_i, n_j) \ge k\)? This is known for
  \(k \le 5\) and open for every \(k \ge 6\); settling \(k = 6\) counts
  as progress.
\end{enumerate}

\emph{(Fine del problema 4. Punteggi: 1+2+3+5+8+13 = 32.)}

\textbf{Cosa consegniamo.} Nessuna prova. Segnalazione rilevante: il
bound asintotico citato nell'enunciato è già stato migliorato in
letteratura.

\subparagraph{2. Dimostrazione}\label{dimostrazione}

Fornal--Sun {[}2607.24655{]} (luglio 2026) dimostrano
\(\max\gcd(m_i,m_j)\gg k\exp(-(2+o(1))\sqrt{\log k/\log\log k})\), con
radice quadrata all'esponente: un miglioramento del fattore esponenziale
richiesto in (b). Riprodurre la loro prova per esteso (grafo dei gcd,
lemma strutturale, partizione crivellante, inversione di Möbius,
trasformata di Fourier discreta) conterebbe secondo le regole; non
fatto.

\subparagraph{3. Verifica: istruzioni, dipendenze,
tempi}\label{verifica-istruzioni-dipendenze-tempi}

Codice disponibile (Python 3, libreria standard; ogni script gira in
meno di un minuto): - Nessun codice ancora.

\subparagraph{4. Fonti e contributo}\label{fonti-e-contributo}

Letteratura arXiv (ricerca deterministica
\texttt{tools/cerca\_letteratura.sh}, abstract letti, non usata come
prova): - arXiv:2603.26043v1 --- Finiteness of Disjoint Covering Systems
with Precisely One Repeated Modulus (Yu Hashimoto, 2026); solo abstract
letto. - arXiv:1511.04293v1 --- Searching for Disjoint Covering Systems
with Precisely One Repeated Modulus (Shalosh B. Ekhad, Aviezri S.
Fraenkel, Doron Zeilberger, 2015); solo abstract letto. -
arXiv:2607.24655v1 --- On the problem of large gcd for disjoint residue
classes (Jan Fornal, Yu-Chen Sun, 2026); solo abstract letto. -
arXiv:2608.15873v1 --- Two Questions on \(G\)-harmonic Tuples (Murali
Menon, 2026); solo abstract letto. {[}2607.24655{]}.

\subparagraph{5. Limiti e parti
irrisolte}\label{limiti-e-parti-irrisolte}

Tutto; la via (b) via riproduzione della prova è quella indicata.

\subsubsection{Decisioni degli agenti}
\paragraph{Run \texttt{p4\_c6}}

\subsubsection{Codice su cui poggia il tentativo}
Nessun codice nei tentativi.

\section{arXiv literature consulted}
\begin{thebibliography}{99}
\bibitem{2603.26043v1} Yu Hashimoto. \emph{Finiteness of Disjoint Covering Systems with Precisely One Repeated Modulus}. arXiv:2603.26043v1 (2026). \url{http://arxiv.org/abs/2603.26043v1}
\bibitem{1511.04293v1} Shalosh B. Ekhad, Aviezri S. Fraenkel, Doron Zeilberger. \emph{Searching for Disjoint Covering Systems with Precisely One Repeated Modulus}. arXiv:1511.04293v1 (2015). \url{http://arxiv.org/abs/1511.04293v1}
\bibitem{2607.24655v1} Jan Fornal, Yu-Chen Sun. \emph{On the problem of large gcd for disjoint residue classes}. arXiv:2607.24655v1 (2026). \url{http://arxiv.org/abs/2607.24655v1}
\bibitem{2608.15873v1} Murali Menon. \emph{Two Questions on \$G\$-harmonic Tuples}. arXiv:2608.15873v1 (2026). \url{http://arxiv.org/abs/2608.15873v1}
\end{thebibliography}
\end{document}
